% chktex-file 44

\section{Parameter Descriptions}

The hardware parameters for the \textbf{Mac} core are described in Table~\ref{table:params}. These parameters control the architectural features, data-path widths, and monitoring capabilities of the controller.

\renewcommand*{\arraystretch}{1.4}
\begin{longtable}[H]{
    | p{0.22\textwidth}
    | p{0.10\textwidth}
    | p{0.08\textwidth}
    | p{0.08\textwidth}
    | p{0.42\textwidth} |
  }
  \hline
  \textbf{Name} & \textbf{Type} & \textbf{Min} & \textbf{Max} & \textbf{Description} \\ \hline \hline
  \endfirsthead
  \hline
  \textbf{Name} & \textbf{Type} & \textbf{Min} & \textbf{Max} & \textbf{Description} \\ \hline \hline
  \endhead

  % --- Hardware Arch ---
  dataW         & Int     & 64      & 64      & XGMII data bus width (Fixed for 10G) \\ \hline
  ctrlW         & Int     & 8       & 8       & XGMII control bus width \\ \hline
  txGbxIfEn     & Boolean & false   & true    & Enable Gearbox interface for TX path \\ \hline
  rxGbxIfEn     & Boolean & false   & true    & Enable Gearbox interface for RX path \\ \hline
  gbxCnt        & Int     & 1       & --      & Width of gearbox synchronization bus \\ \hline
  
  % --- Protocol ---
  paddingEn     & Boolean & false   & true    & Enable auto-padding of short frames \\ \hline
  dicEn         & Boolean & false   & true    & Enable Deficit Idle Count (DIC) for IFG \\ \hline
  minFrameLen   & Int     & 64      & --      & Minimum Ethernet frame size in bytes \\ \hline
  
  % --- PTP ---
  ptpTsEn       & Boolean & false   & true    & Enable Precision Time Protocol (PTP) \\ \hline
  ptpTsFmtTod   & Boolean & false   & true    & Timestamp Format (True: ToD, False: Binary) \\ \hline
  ptpTsW        & Int     & 1       & 128     & Bit-width of the PTP timestamp bus \\ \hline
  
  % --- Stats ---
  statEn        & Boolean & false   & true    & Enable internal Statistics Aggregator \\ \hline
  statTxLevel   & Int     & 0       & 2       & Transmit statistics granularity level \\ \hline
  statRxLevel   & Int     & 0       & 2       & Receive statistics granularity level \\ \hline
  statIdBase    & Int     & 0       & --      & Base index for statistics reporting ID \\ \hline
  statUpdatePeriod & Int  & 1       & --      & Period (cycles) between stats updates \\ \hline
  statStrEn     & Boolean & false   & true    & Enable human-readable metadata strings \\ \hline
  statPrefixStr & String  & --      & --      & Prefix for identifying the MAC instance \\ \hline

  \caption{Mac Parameter Descriptions}\label{table:params}
\end{longtable}

\subsection{Example Instantiation}

The \textbf{Mac} core is typically instantiated in Chisel as a `RawModule` to allow for multiple asynchronous clock domains. An example configuration with PTP enabled is shown below:

\begin{lstlisting}[language=Scala]
  // Define a custom configuration
  val myMacParams = MacParams(
    ptpTsEn      = true,
    ptpTsW       = 96,
    statPrefixStr = "WAN_PORT_0"
  )

  // Instantiate the MAC core
  val ethernetMac = Module(new Mac(myMacParams))
  
  // Connect clocks
  ethernetMac.io.txClk   := coreClock
  ethernetMac.io.rxClk   := phyClock
  ethernetMac.io.statClk := systemClock
\end{lstlisting}